% Generated draft: requires mathematical and source review.
\documentclass[11pt]{article}
\usepackage[margin=25mm]{geometry}
\usepackage{fontspec}
\setmainfont{Latin Modern Roman}
\usepackage{amsmath,amssymb,mathtools,unicode-math}
\setmathfont{Latin Modern Math}
\usepackage{enumitem,needspace,xurl,hyperref}
\setlength{\parindent}{0pt}
\setlength{\parskip}{5pt}
\setlength{\emergencystretch}{3em}
\newcommand{\N}{\mathbb{N}}
\newcommand{\Z}{\mathbb{Z}}
\newcommand{\Q}{\mathbb{Q}}
\newcommand{\R}{\mathbb{R}}
\newcommand{\C}{\mathbb{C}}
\newcommand{\F}{\mathbb{F}}
\newcommand{\A}{\mathbb{A}}
\newcommand{\PP}{\mathbb{P}}
\newcommand{\E}{\mathbb{E}}
\newcommand{\eps}{\varepsilon}
\newcommand{\dd}{\,\mathrm{d}}
\begin{document}
\section*{UnsolvedMath ID: 701687}
\textbf{Source problem code:} catalogue entry
\subsection{Definitions and mathematical statement}
For a finite group \(H\), define its centre and its set of conjugacy class sizes by
\[
Z(H)=\{z\in H:zh=hz\text{ for every }h\in H\},\qquad
\operatorname{cs}(H)=\{|h^H|:h\in H\},
\]
where \(h^H=\{x^{-1}hx:x\in H\}\). Thus \(\operatorname{cs}(H)\) is a set, with multiplicities discarded, and includes the size of the identity class. The group \(H\) is centreless if \(Z(H)=\{1\}\). A non-abelian finite simple group is a finite non-abelian group whose only normal subgroups are the identity subgroup and the whole group. For a positive integer \(n\), let \(S^n\) denote the direct product of \(n\) copies of \(S\).

The question is whether, for every non-abelian finite simple group \(S\), every \(n\in\mathbb N\), and every finite group \(G\),
\[
\bigl(Z(G)=\{1\}\ \text{and}\ \operatorname{cs}(G)=\operatorname{cs}(S^n)\bigr)
\quad\Longrightarrow\quad G\cong S^n.
\]
The source does not specify whether \(0\in\mathbb N\); if it is included, use the empty-product convention \(S^0=\{1\}\).

The supplied record additionally reports that this is proved for \(S=\operatorname{Alt}_5^2\), where \(\operatorname{Alt}_5\) is the alternating group on five letters and the superscript denotes a direct square. This qualification is retained as supplied, but \(\operatorname{Alt}_5^2\) is not simple and therefore does not satisfy the stated hypothesis on \(S\). The intended special case cannot be determined from the record; in particular, it is not silently replaced by the case \(S=\operatorname{Alt}_5\), \(n=2\).
\subsection{Short English statement}
Does the set of conjugacy class sizes of a direct power of a non-abelian finite simple group determine that power up to isomorphism among centreless finite groups? The supplied special-case claim involving \(S=\operatorname{Alt}_5^2\) conflicts with the simplicity hypothesis and requires clarification.
\subsection{Sources}
\begin{itemize}
\item Local UnsolvedMath registry, permanent ID \(701687\), exact source code \texttt{catalogue entry}, title \emph{Kourovka Notebook Problem 20.29}. Supplied upstream revision: \texttt{d371ee917ac9647fbb06cd6fc529deec6f3dda82}. The registry labels the problem \emph{open}; this is a supplied classification, not an independently verified mathematical or literature assessment.
\item \emph{The Kourovka Notebook (Unsolved Problems in Group Theory)}, Problem 20.29, proposed by I. B. Gorshkov in 2022. The supplied background identifies the 20th issue (2022) and the 21st-issue edition (2026), \texttt{arXiv:1401.0300v48}. The supplied collection description credits E. I. Khukhro and V. D. Mazurov as editors.
\item The supplied statement cites \emph{J. Algebra Appl.}, volume \(21\), no. \(11\) (2022), Article ID \(2250226\), 8 pages, for its special-case claim written as \(S=\operatorname{Alt}_5^2\). No article title or authors are supplied. The precise scope of that result has not been independently verified.
\end{itemize}
\end{document}
